\section{Direcciones futuras y problemas abiertos}

\subsection{Dependencias de largo plazo y composición de tareas}

En su forma original, Q-Learning se vuelve ineficiente ante un horizon largo o una sparse reward. Los métodos modernos tienden a combinar actualizaciones Q de múltiples pasos (multi-step TD), una hierarchical policy y un planner externo.

\begin{importantbox}{El pipeline de la composición de tareas}
Dividir una tarea única en dos niveles, Macro y Micro (por ejemplo, la meta policy se encarga de elegir la subtarea y la micro policy, del control concreto), permite atenuar la dificultad de la long-term credit assignment.
\end{importantbox}

\subsection{Gobernanza y transferibilidad}

\begin{knowledgebox}{Lista de problemas abiertos}
\begin{itemize}
    \item Sim-to-real gap: cómo lograr que los datos reales del replay buffer tengan un papel más importante.
    \item Reward hacking: la reward shaping de largo plazo puede introducir atajos (shortcut), por lo que conviene revisar periódicamente la fórmula de la recompensa (reward formula).
    \item Estados multimodales: cuando el robot state incluye imágenes, LiDAR y lenguaje, la complejidad del pipeline de procesamiento de la Q network aumenta de forma abrupta.
\end{itemize}
\end{knowledgebox}

\begin{warningbox}{Gestión del drift durante el despliegue}
Aunque la política tenga un desempeño excelente en la simulación offline, el Observation drift o el action constraint drift pueden provocar fallas en producción. Se recomienda configurar una sanity check periódica (por ejemplo, evaluate on replayed expert rollouts).
\end{warningbox}

\subsection{Resumen de la sección}

Esta sección señala que el campo de Q-Learning sigue en evolución y que las dependencias de largo plazo, la composición de tareas, y la gobernanza y la auditoría son las principales direcciones del trabajo futuro.
